\subsection{Radical de un ideal. Nilradical y radical de un anillo}
\label{subsection:0.1.1-fr}

\begin{env}[1.1.1]
\label{0.1.1.1-fr}
Sea $\mathfrak{a}$ un ideal de un anillo $A$; el \emph{radical} de $\mathfrak{a}$, denotado por $\mathfrak{r}(\mathfrak{a})$, es el conjunto de los $x\in A$ tales que $x^n\in\mathfrak{a}$ para al menos un entero $n>0$; es un ideal que contiene $\mathfrak{a}$.
Se tiene $\mathfrak{r}(\mathfrak{r}(\mathfrak{a}))=\mathfrak{r}(\mathfrak{a})$; la relación $\mathfrak{a}\subset\mathfrak{b}$ implica $\mathfrak{r}(\mathfrak{a})\subset\mathfrak{r}(\mathfrak{b})$; el radical de una intersección finita de ideales es la intersección de sus radicales.
Si $\varphi$ es un homomorfismo de un anillo $A'$ en $A$, se tiene $\mathfrak{r}(\varphi^{-1}(\mathfrak{a}))=\varphi^{-1}(\mathfrak{r}(\mathfrak{a}))$ para todo ideal $\mathfrak{a}\subset A$.
Para que un ideal sea el radical de un ideal, es necesario y suficiente que sea una intersección de ideales primos.
El radical de un ideal $\mathfrak{a}$ es la intersección de los ideales primos \emph{minimales} entre los que contienen $\mathfrak{a}$; si $A$ es noetheriano, estos ideales primos minimales son finitos en número.

El radical del ideal $(0)$ también se llama \emph{nilradical} de $A$; es el conjunto $\mathfrak{N}$ de los elementos nilpotentes de $A$.
Se dice que el anillo $A$ es \emph{reducido} si $\mathfrak{N}=(0)$; para todo anillo $A$, el cociente $A/\mathfrak{N}$ de $A$ por su nilradical es un anillo reducido.
\end{env}

\begin{env}[1.1.2]
\label{0.1.1.2-fr}
Recordemos que el \emph{radical} $\mathfrak{R}(A)$ de un anillo $A$ (no necesariamente conmutativo) es la intersección de los ideales maximales a izquierda de $A$ (y también la intersección de los ideales maximales a derecha).
El radical de $A/\mathfrak{R}(A)$ es $(0)$.
\end{env}
